% Labb 3, parts 11, 12 and 14 to 17, from labs/lab3/c_program_flow.md. Each part is a \labtask,
% and its Mål, Bakgrund, Specifikation, Programmet, Uppgifter and Frågor are \task headings. The
% program given in part 15 is the handout's.

\section{Del 11, 12 och 14--17: Programflöde, loopar och tidsfördröjningar}
\label{lab3:sec:c}
Delarna hör till kapitel~\ref{c9:ch}. Varje del är ett eget projekt som utgår från mallen
(\secref{lab3:template}).

Från och med nu gäller två vanor till:
\begin{itemize}
  \item \textbf{Rita flödesplanen först}, på papper, innan ni skriver en enda instruktion
    (\secref{c9:a1}).
  \item \textbf{Mät tider med stoppuret}, med Frequency inställd på 16~MHz (\secref{c9:b6}).
\end{itemize}

Del 13 ingår inte i laborationen; numreringen hoppar från 12 till 14.

\labtask{Del 11}{Vitsen med flaggor: jämförelse med cp}\phantomsection\label{lab3:d11}
\task{Mål}
Se vad \code{cp} och \code{cpi} gör med flaggorna, och hur flaggorna säger om två tal är lika,
större eller mindre.

\task{Bakgrund}
\Secref{c9:a3} och~\ref{c9:a6}.

\task{Uppgifter}
Skriv ett program med de fem jämförelserna nedan, med en etikett efter var och en.
\textbf{Förutsäg} \code{Z}, \code{C} och \code{S} för varje jämförelse i en tabell innan ni kör
fram till etiketterna.

\begin{filltable}{@{}p{8.8em}cc*{3}{C{1.8em}}Y@{}}
  \thead{Jämförelse} & \thead{\code{r16}} & \thead{\code{r17}} & \thead{Z} & \thead{C} &
    \thead{S} & \thead{\code{r16} är \dots\ \code{r17}} \\
  \midrule
  1. \code{cp r16, r17} & 50 & 30 & & & & \fillrule
  2. \code{cp r16, r17} & 50 & 50 & & & & \fillrule
  3. \code{cp r16, r17} & 50 & 70 & & & & \fillrule
  4. \code{cpi r16, 50} & 50 & -- & & & & \fillrule
  5. \code{cp r16, r17} & $-5$ & 3 & & & & \\
\end{filltable}

\begin{enumerate}
  \item Fyll i tabellen, och fyll i sista kolumnen med \emph{större än}, \emph{lika med} eller
    \emph{mindre än}.
  \item Kör fram till varje etikett och kontrollera flaggorna.
  \item Kontrollera efter jämförelse 1--3 att \code{r16} fortfarande är 50. Varför ändras det inte,
    fast \code{cp} räknar ut en subtraktion?
  \item Titta särskilt på jämförelse~5. Är $-5$ mindre eller större än 3? Svara både
    \textbf{utan} tecken, när \code{r16} läses som 251, och \textbf{med} tecken. Vilken flagga
    svarar på vilken fråga?
\end{enumerate}

\task{Frågor}
\begin{itemize}
  \item Vilka två flaggor räcker för att avgöra om två tal utan tecken är lika, om det första är
    större, eller om det är mindre?
  \item Varför behövs både \code{brlo} och \code{brlt}?
\end{itemize}

\redovisa{tabellen.}

\labtask{Del 12}{Flödesplan och villkorliga hopp}\phantomsection\label{lab3:d12}
\task{Mål}
Rita en flödesplan med beslut, och skriva det som ett program med jämförelser och villkorliga hopp.

\task{Bakgrund}
\Secref{c9:a4} och~\ref{c9:a5}.

\task{Specifikation}
En termostat styr ett rum. Temperaturen i grader finns i \code{r16}. Port B styr två saker:
\begin{itemize}
  \item \textbf{PB0}, en värmare, ska vara på när temperaturen är \textbf{under 20} grader;
  \item \textbf{PB1}, en fläkt, ska vara på när temperaturen är \textbf{25 grader eller mer};
  \item mellan 20 och 24 grader ska båda vara av.
\end{itemize}

Programmet ska gå i en loop: det läser \code{r16}, sätter utgångarna, och börjar om. Då kan ni
stoppa simuleringen, ändra värdet i \code{r16} i Processor Status, och köra vidare för att se
utgångarna ändras.

\task{Uppgifter}
\begin{enumerate}
  \item Rita flödesplanen. Den har två beslut. Märk varje beslut med Ja och Nej.
  \item Skriv programmet. Ladda startvärdet 18 i \code{r16} \textbf{före} loopen, så att ett värde
    ni skriver in i Processor Status inte skrivs över.
  \item \textbf{Förutsäg} \code{PORTB} för temperaturerna 18, 20, 24, 25 och 30. Testa alla fem i
    simulatorn.
  \item Vilka av värdena är \textbf{gränsfall}, och varför är det de som är viktigast att testa?
\end{enumerate}

\task{Frågor}
\begin{itemize}
  \item Villkoret \enquote{25 grader eller mer} används för fläkten. Hur skriver man \enquote{mer
    än 24} med de hopp som finns?
  \item Vad händer med termostaten vid $-5$ grader? Se \secref{c9:a6}.
\end{itemize}

\redovisa{flödesplanen och programmet, och visa minst tre temperaturer i simulatorn.}

\labtask{Del 14}{Loop med visst antal varv}\phantomsection\label{lab3:d14}
\task{Mål}
Skriva en loop som går ett bestämt antal varv med en räknare, och veta exakt hur många gånger
varje instruktion i den körs.

\task{Bakgrund}
\Secref{c9:b1}.

\task{Specifikation}
Processorn har ingen multiplikation i den här laborationen. Räkna ut $7 \cdot 3$ genom att addera 3
sju gånger, med en loop. Visa antalet varv som hittills gått på port B, så att det syns i
I/O-fönstret.

\task{Uppgifter}
\begin{enumerate}
  \item Rita flödesplanen. Använd ett register för resultatet, ett för talet som adderas, ett för
    varven som är kvar och ett för varven som har gått.
  \item Skriv programmet med \code{.equ} för konstanterna, till exempel \code{.equ LAPS = 7} och
    \code{.equ TERM = 3}.
  \item \textbf{Förutsäg} resultatet och \code{PORTB} när programmet är klart. Stega igenom minst
    två varv och se \code{PORTB} räkna, kör sedan till slutet med \textbf{Run To Cursor}.
  \item Hur många gånger körs \code{brne}, och hur många av dem hoppar den?
  \item Ändra \code{LAPS} till 0. \textbf{Förutsäg} resultatet innan ni kör. Förklara det ni ser.
  \item Räkna ut $12 \cdot 25$ med samma program. Stämmer resultatet? Varför, eller varför inte?
\end{enumerate}

\task{Frågor}
\begin{itemize}
  \item Varför behövs ingen \code{cpi} före \code{brne}?
\end{itemize}

\redovisa{programmet och svaren på uppgift 4--6.}

\labtask{Del 15}{Kort tidsfördröjning}\phantomsection\label{lab3:d15}
\task{Mål}
Göra en tidsfördröjning med en loop, räkna ut dess längd i klockcykler, och kontrollera den med
stoppuret.

\task{Bakgrund}
\Secref{c9:b3} och~\ref{c9:b4}.

\task{Programmet}
\begin{asmcode}
main:
    sbi DDRB, 0                     ; PB0 is an output.
    sbi PORTB, 0                    ; The pulse starts.

delay_start:
    ldi r18, 200
delay_loop:
    dec r18
    brne delay_loop
delay_end:
    cbi PORTB, 0                    ; The pulse ends.

end:
    rjmp end
\end{asmcode}

Programmet ger en puls på PB0: stiftet går högt, fördröjningen körs, och stiftet går lågt.

\task{Uppgifter}
\begin{enumerate}
  \item Räkna ut fördröjningens längd, från \code{delay_start} till \code{delay_end}, i
    klockcykler och i mikrosekunder. Använd en tabell med en rad per instruktion, som i
    \secref{c9:b4}.
  \item Mät samma sak i simulatorn: brytpunkt på \code{delay_start} och \code{delay_end},
    Frequency på 16~MHz, och stoppuret nollställt vid den första brytpunkten. Stämmer det?
  \item Pulsen ska vara \textbf{30~µs}. Räkna ut hur många varv loopen ska gå, ändra programmet, och
    mät.
  \item Vilken är den längsta fördröjning loopen kan ge? Räkna först, mät sedan.
  \item Flytta den första brytpunkten till \code{sbi PORTB, 0}. \textbf{Förutsäg} hur mycket
    mätningen ändras.
\end{enumerate}

\task{Frågor}
\begin{itemize}
  \item Varför blir det exakt $3 \cdot n$ cykler och inte $3 \cdot n + 1$?
\end{itemize}

\redovisa{uträkningen och mätningen för 30~µs.}

\labtask{Del 16}{Längre tidsfördröjning}\phantomsection\label{lab3:d16}
\task{Mål}
Göra en längre fördröjning med två loopar i varandra, och träffa en önskad tid exakt.

\task{Bakgrund}
\Secref{c9:b5}.

\task{Specifikation}
PB0 ska byta läge var \textbf{10:e millisekund}: hög i 10~ms, låg i 10~ms, och så vidare, för
alltid. Byt läge genom att läsa \code{PORTB}, invertera bit~0 med \code{eor} och en mask, och
skriva tillbaka. Fördröjningen mellan bytena görs med två loopar i varandra.

\task{Uppgifter}
\begin{enumerate}
  \item Rita flödesplanen: en loop som byter läge på PB0 och sedan väntar, med fördröjningen som
    två loopar i varandra.
  \item Välj OUTER och INNER med formeln $\text{OUTER} \cdot (3 \cdot \text{INNER} + 3)$ så att
    fördröjningen blir så nära 10~ms som möjligt. Skriv programmet och mät fördröjningen.
  \item Lägg till en \code{nop} i den yttre loopen. Hur ändras formeln? Välj nya konstanter så att
    fördröjningen blir \textbf{exakt} 10~ms, 160\,000 cykler. Mät.
  \item Mät nu en hel halvperiod: från ett byte av PB0 till nästa. Den blir några cykler längre än
    fördröjningen. Hur många, och vilka instruktioner är det?
  \item Kör programmet fritt med \textbf{F5} i några sekunder och stoppa. Hur lång tid har gått
    enligt stoppuret? Ungefär hur många gånger har PB0 bytt läge?
\end{enumerate}

\task{Frågor}
\begin{itemize}
  \item Skulle en lysdiod på PB0 synas blinka med den här fördröjningen på ett riktigt kort?
    Varför inte?
\end{itemize}

\redovisa{konstanterna för exakt 10~ms och mätningarna.}

\labtask{Del 17}{Flervalssituationer och omarbetning av flödesplaner}%
\phantomsection\label{lab3:d17}
\task{Mål}
Skriva ett program som väljer mellan flera alternativ, och arbeta om en flödesplan så att
programmet får en väg in och en väg ut.

\task{Bakgrund}
\Secref{c9:a7} och~\ref{c9:a8}.

\task{Specifikation}
Ett läge i \code{r16} väljer ett mönster på port B:

\begin{narrowtable}{@{}cl@{}}
  \thead{Läge} & \thead{Mönster på port B} \\
  \midrule
  0 & alla släckta \\
  1 & alla tända \\
  2 & de fyra nedre, PB0--PB3 \\
  3 & de fyra övre, PB4--PB7 \\
  annat & varannan, \code{0101 0101}, som felindikering \\
\end{narrowtable}

Programmet går i en loop, så att läget kan ändras i Processor Status medan programmet står
stilla.

\task{Uppgifter}
\begin{enumerate}
  \item Rita en \textbf{första} flödesplan, hur ni vill. Skriv programmet efter den, och testa alla
    lägen, inklusive ett felläge.
  \item Räkna i er flödesplan och i programmet: hur många rutor eller instruktioner skriver till
    port B? Hur många ställen går programmet tillbaka till loopens början från?
  \item \textbf{Arbeta om} flödesplanen enligt \secref{c9:a8}: låt varje alternativ bara välja
    mönster, och låt alla alternativ mötas i \textbf{en} ruta som skriver till port B.
  \item Skriv om programmet efter den nya flödesplanen. Testa alla lägen igen. Programmet ska göra
    exakt samma sak som förut.
  \item Lägg till läge~4, PB0 och PB7 tända. Hur många ställen i programmet behövde ändras, i den
    första versionen och i den omarbetade?
\end{enumerate}

\task{Frågor}
\begin{itemize}
  \item Varför är det viktigt att ha med alternativet \enquote{annat} i flödesplanen?
\end{itemize}

\redovisa{båda flödesplanerna och det omarbetade programmet.}
